Sparse attention in multi-agent communication is usually justified by the
operations it removes. We charge each peer-selection rule for what deciding
costs, and measure on four GPUs, from A100 to Blackwell, when the saving becomes
elapsed time.
First, a rule that ranks attention scores must compute all of them, so it
removes at most half of the forward-pass multiply--accumulates at any sparsity;
rules that select before scoring escape this ceiling by as much as their
selector is cheap. Second, composed from library operators, score-then-select
attention is slower than dense attention at every configuration we measure, by
up to 37 times, because compacting the retained peers adds more memory traffic
than it removes. A fused kernel that streams the selection is faster than dense
attention with one or two retained peers in fp32 and, up to 512 agents, within
15\% of parity with four and slower with eight; and from 384 agents upward at least 85\% of its time is spent
finding the retained set; in bf16 only one retained peer breaks even. On A100,
H100 and H200 the streaming kernel wins everywhere only with one retained peer,
and in bf16 not at all. Third,
once the set is known without scoring, attention over it is up to 16 times
faster than dense attention, and the selector decides the outcome: of three
selectors, only the cheapest reaches parity, at 512 agents, on best-case
timings and on the Blackwell card only. Below 100 agents, where multi-agent methods are usually evaluated, the
number of kernels a step launches orders latency, and its operation count does
not.
